\section{File Storage and Node Input/Output Management}
\label{sec:file-storage}

Many workflow nodes operate on binary content rather than plain text — image-processing nodes accept and return images, AI nodes accept photographs, and users often attach files to a node before a workflow ever runs. Since the execution engine itself only passes strings and JSON between nodes, a separate subsystem is needed to accept uploads, persist them durably, and expose stable URLs that nodes can consume during execution. This subsystem uses Cloudinary as the storage backend, with two parallel Postgres models: \texttt{NodeInput} for user-attached files, and \texttt{NodeOutput} for files produced by nodes during execution.

\subsection{Design Goals}

\begin{itemize}
  \item \textbf{Decoupling execution from storage.} A node such as \texttt{blurDenoise} receives a URL string, not a raw buffer — keeping the execution engine's data model uniform while still allowing large binary payloads to be referenced cheaply.
  \item \textbf{Draft vs.\ committed files.} A user may upload a file to a node while still editing a workflow, then discard the node before saving. Files are not considered "attached" until the workflow is saved, hence an explicit \texttt{isDraft} flag.
  \item \textbf{Per-execution provenance.} Every file a node produces must be traceable to the exact node, execution, project, and user that generated it, for both the UI and ownership/cleanup enforcement.
\end{itemize}

\subsection{The Cloudinary Service Layer}

All Cloudinary interaction is isolated in \texttt{cloudinary.service.ts}, so no controller or node talks to the SDK directly — giving one place to swap providers or reason about error handling.

\begin{lstlisting}
export const uploadToCloudinary = async (folder: string, fileBuffer: Buffer): Promise<UploadResult> => {
  return new Promise((resolve, reject) => {
    const stream = cloudinary.uploader.upload_stream(
      { folder, resource_type: "auto", use_filename: true, unique_filename: false, overwrite: true },
      (error, result) => {
        if (error) reject(new AppError(`Upload failed: ${error.message}`, 500));
        else resolve({ url: result!.secure_url, publicId: result!.public_id });
      }
    );
    stream.end(fileBuffer);
  });
};
\end{lstlisting}

Key decisions:
\begin{itemize}
  \item \textbf{Streaming, not disk-buffering.} The in-memory buffer is piped directly into \texttt{upload\_stream}, so files never touch local disk — important for stateless horizontal scaling.
  \item \textbf{\texttt{resource\_type: "auto"}.} Lets Cloudinary infer the type, keeping the service generic across current and future node categories.
  \item \textbf{Deterministic, caller-supplied folder paths.} Paths encode ownership (e.g.\ \texttt{users/\{userId\}/projects/\{projectId\}/inputs/\{nodeId\}/\{uuid\}}), making bulk cleanup by prefix trivial — deleting \texttt{users/\{userId\}} removes everything that user ever uploaded or generated.
\end{itemize}

Deletion requires the resource type explicitly (\texttt{"auto"} is upload-only), so it's derived from the stored MIME type:

\begin{lstlisting}
const toCloudinaryResourceType = (mimeType: string): "image" | "video" | "raw" => {
  if (mimeType.startsWith("image/")) return "image";
  if (mimeType.startsWith("video/")) return "video";
  return "raw";
};
\end{lstlisting}

This is why every \texttt{NodeInput}/\texttt{NodeOutput} row persists its own \texttt{mimeType} — deletion often happens long after the original request's \texttt{Content-Type} context is gone.

\subsection{Input Flow: User-Attached Files}
\label{subsec:input-flow}

Uploads pass through \texttt{uploadMiddleware}, a Multer wrapper using in-memory storage with a 5\,MB per-file limit and up to 50 files per request, supporting nodes whose input declares \texttt{multiple: true}. Multer errors are converted to \texttt{AppError} so they flow through the same centralized error pipeline as everything else.

The \texttt{uploadNodeFiles} controller verifies project ownership, generates a collision-proof path per file, uploads via the Cloudinary service, and writes one \texttt{NodeInput} row per file with \texttt{isDraft: true}:

\begin{lstlisting}
const uniquePublicId = `users/${userId}/projects/${projectId}/inputs/${nodeId}/${crypto.randomUUID()}`;
const uploadResult = await uploadToCloudinary(uniquePublicId, file.buffer);
await prisma.nodeInput.create({
  data: { nodeId, projectId, userId, mimeType: file.mimetype,
          publicId: uploadResult.publicId, url: uploadResult.url, isDraft: true }
});
\end{lstlisting}

A node in the visual builder can have experimental uploads that never get promoted if the workflow is never saved — those rows remain orphaned drafts. \texttt{getNodeFiles} filters strictly on \texttt{isDraft: false}, so drafts stay invisible until promoted.

\textbf{Promotion on save} happens inside \texttt{updateWorkflow}:

\begin{lstlisting}
const nodeIds = workflow.nodes.map((n: any) => n.id);
await prisma.nodeInput.updateMany({
  where: { nodeId: { in: nodeIds }, userId, isDraft: true },
  data: { isDraft: false },
});
\end{lstlisting}

Because the update is scoped to node IDs present in the saved workflow, draft files belonging to nodes removed before saving are not promoted and remain orphaned — reclaimed later by prefix-based Cloudinary cleanup regardless of draft status.

\textbf{Retrieval and deletion} endpoints:
\begin{itemize}
  \item \texttt{GET /node/:nodeId/files} — committed files, scoped to project and user.
  \item \texttt{DELETE /node/:nodeId/files} — bulk-deletes all files for a node; Cloudinary deletion happens first, and Postgres rows are only removed once all remote deletions succeed, avoiding orphaned-in-DB-but-missing-in-storage states.
  \item \texttt{DELETE /node/:nodeId/files/:fileId} — same ordering, single file.
\end{itemize}

An admin-only \texttt{GET /node/files/all} exposes every \texttt{NodeInput} row platform-wide (paginated, optional \texttt{isDraft} filter, joined with owner email) for storage auditing.

\subsection{Output Flow: Node-Produced Files}
\label{subsec:output-flow}

While inputs are driven by HTTP requests from the browser, outputs are produced inside the BullMQ worker as a node's \texttt{execute()} runs. Three patterns converge on \texttt{NodeOutput}.

\textbf{Generic upload helper.} \texttt{uploadNodeOutputs} in \texttt{nodeoutput.service.ts} is for nodes already holding a raw buffer. It validates file count/size and execution existence, then uploads under \texttt{outputs/execution/\{executionId\}/\{nodeId\}} — namespaced by execution so repeated runs of the same node (e.g.\ in a loop) stay distinguishable. \texttt{userId} and \texttt{projectId} are derived server-side from \texttt{executionId} rather than trusted from the caller, preventing a buggy node from writing under the wrong owner.

\textbf{Direct upload from external-API nodes.} Most image-processing nodes (\texttt{blurDenoise}, \texttt{edgeDetection}, \texttt{resize}, \texttt{colorMasking}, \texttt{morphology}, \texttt{threshold}, \texttt{contour}, \texttt{colorSpaceConversion}, \texttt{normalize}, \texttt{crop}) don't upload buffers themselves — they call a Python image-processing microservice with the source image URL plus \texttt{user\_id}/\texttt{project\_id}/\texttt{execution\_id}/\texttt{node\_id}; that service transforms the image, uploads it to Cloudinary, and returns only \texttt{output\_url}. The Node.js side just records the metadata in \texttt{NodeOutput}.

\textbf{AI-node pattern.} \texttt{objectDetection} and \texttt{singleImageClassification} fetch the source image bytes themselves, forward them to a Hugging Face endpoint, and (for \texttt{objectDetection}, which returns an annotated image) re-upload via the same \texttt{uploadToCloudinary} helper before recording a \texttt{NodeOutput} row.

Regardless of which pattern a node uses, \texttt{execute()} always resolves to an array of Cloudinary URLs or JSON strings handed to the next node by the worker — keeping the contract uniform from the execution engine's perspective.

\subsection{Schema Compatibility}

Full ERD rationale is in Subsection~\ref{subsec:data-layer-architecture}; briefly, the existing schema needs no changes to support this feature:

\begin{itemize}
  \item Splitting \texttt{NodeInput}/\texttt{NodeOutput} into distinct models mirrors the two independent flows (pre-execution attachment vs.\ post-execution artifact), avoiding a discriminator column and letting each evolve its own indexing.
  \item \texttt{NodeInput.isDraft} directly supports promotion-on-save without extra schema.
  \item \texttt{NodeOutput.executionId} (absent from \texttt{NodeInput}) keeps outputs from repeated runs distinguishable, matching the Cloudinary folder namespacing.
  \item Both models redundantly store \texttt{userId}/\texttt{projectId} so ownership checks and bulk deletion can filter directly without traversing the \texttt{Project.workflow} JSON blob.
  \item \texttt{publicId} is unique on both models, serving as the natural key reconciling a Postgres row with its Cloudinary asset during deletion.
\end{itemize}

\subsection{Ownership, Cleanup, and Lifecycle}
\label{subsec:lifecycle}

Cleanup is anchored on the Cloudinary folder convention rather than iterating individual rows. On account deletion, a single prefix delete is issued:

\begin{lstlisting}
await cloudinary.api.delete_resources_by_prefix(`users/${id}`);
\end{lstlisting}

This runs \emph{outside} the surrounding Prisma transaction and is deliberately non-fatal on failure — a failed cloud cleanup must never block the database anonymization from completing, since scrubbing relational PII is the higher-priority guarantee. The corresponding \texttt{NodeInput}/\texttt{NodeOutput} rows are removed as part of the same atomic transaction that deletes the user's projects and credential sets, so no orphaned metadata survives even if the Cloudinary side only partially succeeded.